% Summer research report: two-particle Hamiltonian system
% Simplified version (2026-11-07). The previous, longer version is in
% report/archive/report_2026-11-07_before_simplify.tex.
% Data sources: report01-04.ipynb, case3_newton_failure_diagnosis.ipynb,
% lyapunov_*.py; figures from the scripts listed in Table 2.
% Text marked with \todo{...} still needs checking or a decision.

\documentclass[11pt,a4paper]{article}

\usepackage[margin=2.5cm]{geometry}
\usepackage{amsmath,amssymb}
\usepackage{graphicx}
\usepackage{booktabs}
\usepackage{xcolor}
\usepackage[colorlinks=true,linkcolor=blue!50!black,citecolor=blue!50!black,urlcolor=blue!50!black]{hyperref}

\graphicspath{{../images/}{figures/}}

% Figures and tables use [H] (float package): placed exactly where they are
% in the source, right after the paragraph that first refers to them.
\usepackage{float}

% Visible placeholder. For a clean draft use \newcommand{\todo}[1]{}
\newcommand{\todo}[1]{\textcolor{red}{[TODO: #1]}}

\newcommand{\q}{\mathbf{q}}
\newcommand{\p}{\mathbf{p}}

\title{Periodic Orbits and Chaos in a System of Two Repelling Particles in a Box:\\
a Numerical Study}
\author{Jiajun Wu\\[2pt]
\small Department of Mathematics, Imperial College London\\
\small Summer research project (2026), supervised by Prof.\ Dmitry Turaev}
\date{November 2026}

\begin{document}
\maketitle

% ---------------------------------------------------------------------------
\begin{abstract}
We study numerically two particles in a square box that repel each other and
are pushed back by the walls, restricting to motions that are symmetric about
the centre of the box. Using Poincar\'e sections, we find that the motion is
regular near the equilibrium and chaotic at slightly higher energies. At high
energy chaotic motion is still present, but stable periodic orbits are
surrounded by islands of regular motion, and trajectories that start in these
islands stay close to the periodic orbit. We compute two families of periodic orbits with Newton's method and
follow them as the energy changes. A ``vertical'' orbit, in which each particle
bounces between the top and bottom walls, is linearly stable for $E\gtrsim3.9$
and unstable at lower energies. A ``circular'' orbit around the centre of the box
is linearly stable at every energy we computed, from $E\approx2$ to
$E\approx202$. At $E\approx52$, Lyapunov exponents show a chaotic region with
many such islands.
\end{abstract}

% ---------------------------------------------------------------------------
\section{Introduction}\label{sec:intro}

Two particles that repel each other inside a box form a very simple system,
but their long-time motion is not obvious. They might settle into a repeating
pattern, or move chaotically over the whole box, depending on the energy and
the initial condition.

Rom-Kedar and Turaev \cite{turaev2022} proved that, under mild conditions,
the motion of $N$ repelling particles in a bounded region is not ergodic at
sufficiently high energies. They construct stable periodic orbits that are
surrounded by a set of positive measure of quasi-periodic motions; trajectories
that start in this set stay near the periodic orbit for all time, so they
cannot explore the whole energy surface. For particles in a rectangular box,
the stable orbit they construct is one in which every particle moves fast in
the same direction, say vertically, each on its own line and in synchrony with
the others \cite[Sec.~2.6]{turaev2022}. Their result is asymptotic: it holds
only above some energy that is not given explicitly, and it needs the wall
potential to be steep enough (Section~\ref{sec:vertical}). It therefore says
nothing about moderate energies.

In this project we study numerically the simplest case, two particles in a
square box, restricted to motions that are symmetric about the centre. We look
at how the type of motion, regular or chaotic, changes with energy, from the
equilibrium up to high energies. We then compute two families of periodic
orbits, a vertical orbit of the type constructed in \cite{turaev2022} and a
circular orbit around the centre, and follow how their shape and stability
vary with energy.

The report is organised as follows. Section~\ref{sec:model} describes the
model and Section~\ref{sec:methods} the numerical methods.
Section~\ref{sec:results} presents the results: typical trajectories, how the
motion changes with energy, two families of periodic orbits, and a picture of
the phase space at high energy. Section~\ref{sec:discussion} summarises the
findings and their limitations. All results are numerical; none of them is
proved.

% ---------------------------------------------------------------------------
\section{The model}\label{sec:model}

Two particles with positions $\q_i=(x_i,y_i)$ and momenta $\p_i$ ($i=1,2$)
move in the box $(0,L_x)\times(0,L_y)$. The Hamiltonian is
\begin{equation}\label{eq:H}
  H(\q,\p)=\sum_{i=1}^{2}\frac{|\p_i|^2}{2m_i}
  + V\bigl(|\q_1-\q_2|\bigr) + U(\q_1) + U(\q_2),
\end{equation}
where
\begin{equation}\label{eq:potentials}
  V(r)=k\,r^{-\alpha},\qquad
  U(x,y)=C\Bigl[x^{-\beta}+(L_x-x)^{-\beta}+y^{-\beta}+(L_y-y)^{-\beta}\Bigr].
\end{equation}
$V$ is the repulsion between the particles and $U$ pushes each particle away
from the walls. Both become infinite at collision or at a wall, so the
particles can never collide or leave the box.

We take $k=C=\alpha=\beta=1$, $m_1=m_2=1$ and $L_x=L_y=10$. These are simply
convenient values of order one; the model is not meant to describe a specific
physical system.

\paragraph{Energy.} The energy $H$ is conserved. Because of the walls,
momentum and angular momentum are not.

\paragraph{Symmetric motions.}
If particle~2 starts at the point reflection of particle~1 through the centre
of the box, $\q_2=(10,10)-\q_1$, with the opposite momentum $\p_2=-\p_1$, then
by symmetry the two particles stay in this position relative to each other for
all time. We call such motions \emph{centre-symmetric}. For them it is enough
to follow particle~1, so a Poincar\'e section can be drawn as a
two-dimensional picture (Section~\ref{sec:lyapunov}). All initial conditions
in this report are centre-symmetric.

\paragraph{Equilibrium.}
The particles are at rest in equilibrium at $\q_1=(a,a)$, $\q_2=(10-a,10-a)$
with $a\approx3.6645$, on the diagonal of the box. Its energy is
$E_{\mathrm{eq}}\approx1.9877$, the lowest energy in this report.

% ---------------------------------------------------------------------------
\section{Numerical methods}\label{sec:methods}

\subsection{Time integration}
We integrate the equations of motion with the leapfrog method:
\begin{equation}\label{eq:leapfrog}
  \p^{n+1/2}=\p^n+\tfrac h2 F(\q^n),\quad
  \q^{n+1}=\q^n+h\,\p^{n+1/2},\quad
  \p^{n+1}=\p^{n+1/2}+\tfrac h2 F(\q^{n+1}),
\end{equation}
where $F=-\nabla(V+U)$ is the force and $h$ the time step. Leapfrog is a
second-order method whose energy error stays small and does not grow over long
times \cite{hairer2006}, which is important because we integrate for hundreds
of time units. We use smaller steps when the particles move faster, from
$h=10^{-2}$ near the equilibrium down to $h=10^{-5}$ at $E\approx202$. The
relative energy error in the runs shown is at most about $10^{-6}$, and
halving $h$ changes the computed periods by less than one part in $10^{7}$.
% Step-halving check run 2026-11-07: E=11.075 T=6.87461211/...275/...291;
% E=3.876 T=14.80128691/...697/...699; E=3.011 T=23.80233268/...379/...407.

\subsection{Poincar\'e sections}
A Poincar\'e section records the state each time the orbit crosses a chosen
surface in a chosen direction, for example $y_1=5$ with $p_{y_1}>0$. The full
phase space is 8-dimensional; fixing the energy and the surface leaves a
6-dimensional section, which is the one Newton's method below works on. For a
centre-symmetric motion, however, the state of particle~2 is fixed by that of
particle~1, so the motion lies in the 4-dimensional space
$(x_1,y_1,p_{x_1},p_{y_1})$. On a fixed energy level and the surface $y_1=c$,
$p_{y_1}>0$, a point is then determined by $(x_1,p_{x_1})$, because $p_{y_1}$
follows from the energy. The section of a centre-symmetric motion is therefore
exactly two-dimensional, not a projection. In the same way, the surface
$x_1=c$, $p_{x_1}>0$ gives the coordinates $(y_1,p_{y_1})$.

On a two-dimensional section the pattern of the points tells us the type of
motion \cite{lichtenberg1992}: a periodic orbit that crosses the section $k$
times per period gives $k$ points, a regular (quasi-periodic) orbit gives
points on one or several smooth closed curves, and a chaotic orbit gives points
scattered over a two-dimensional region. This reading relies on the section
being two-dimensional: in a projection of a higher-dimensional section,
quasi-periodic motion could also look scattered.

A time step usually jumps over the surface, so the crossing point has to be
estimated. For plotting sections we use linear interpolation between the two
steps on either side of the surface; its error is of the same order, $h^2$, as
leapfrog itself. For Newton's method (below) we need the crossing more
accurately, so there we instead find the length $\tau<h$ of a shortened
leapfrog step that lands exactly on the surface, by bisection (or, in
Section~\ref{sec:lyapunov}, by the secant method).

\subsection{Finding periodic orbits}\label{sec:newton}
Let $P$ be the Poincar\'e map, which sends a point $z$ on the section to the
point where the orbit next crosses it. At a fixed energy, $z$ has six
coordinates. A periodic orbit is a fixed point, $P(z^*)=z^*$ (or
$P^k(z^*)=z^*$ if the orbit crosses the section $k$ times per period). We solve
this equation with Newton's method,
\begin{equation}\label{eq:newton}
  z_{j+1}=z_j-\bigl(DP(z_j)-I\bigr)^{-1}\bigl(P(z_j)-z_j\bigr),
\end{equation}
where the Jacobian $DP$ is computed by finite differences. If a full Newton
step does not reduce the error, we take a shorter step. From a good initial
guess Newton converges in 3--6 iterations, with $|P(z)-z|$ between $10^{-9}$
and $10^{-13}$.

\paragraph{Stability.}
The eigenvalues $\lambda$ of $DP(z^*)$, called Floquet multipliers, tell us
what happens to a small perturbation of the periodic orbit after each period:
it is multiplied by roughly $|\lambda|$ in the corresponding direction. If all
multipliers satisfy $|\lambda|=1$, small perturbations neither grow nor shrink
and the orbit is \emph{linearly stable}. If some $|\lambda|>1$, perturbations
grow and the orbit is \emph{unstable}. Since our Jacobians are accurate to
about $10^{-6}$, we count $|\lambda|$ within $10^{-3}$ of $1$ as equal to $1$.

\paragraph{Following an orbit in energy.}
To see how a periodic orbit changes with energy, we change $E$ in small steps
and use the orbit found at the previous energy, with its momentum rescaled, as
the starting guess for Newton at the new energy.

\subsection{Lyapunov exponents}
The Lyapunov exponent $\lambda_{\max}$ of a trajectory measures how fast two
nearby trajectories separate: it is close to $0$ for regular motion and
positive for chaotic motion. We estimate it by following, along each
trajectory, a small perturbation for a time $T=500$ and measuring its average
exponential growth rate. We do this for a grid of 13\,920 centre-symmetric
initial conditions and colour each point in the Poincar\'e section by the
exponent of its trajectory.

% ---------------------------------------------------------------------------
\section{Results}\label{sec:results}

\subsection{Five typical trajectories}\label{sec:trajectories}
Table~\ref{tab:cases} lists five initial conditions, and
Figure~\ref{fig:cases} shows their trajectories and Poincar\'e sections.
\begin{itemize}
  \item Cases~1 and~2 start near the equilibrium and stay near it. Their
        section points lie on smooth curves, so the motion is regular.
  \item In Case~3 the particles go round the centre of the box, always
        opposite each other. The section points form two tiny clusters, which
        suggests a periodic orbit nearby (found in Section~\ref{sec:circular}).
  \item In Case~4 each particle bounces between the top and bottom walls along
        an almost vertical line. This is a periodic orbit: after one period
        $T\approx1.0027$ the state returns to its starting value to within
        $6\times10^{-7}$.
  \item Case~5 starts like Case~4 but at a different $x$. Its section points
        lie on a closed curve around the Case~4 point: the orbit oscillates
        around the stable periodic orbit of Case~4.
\end{itemize}

\begin{table}[H]
\centering
\caption{Initial conditions. Particle~2 always starts at the point reflection
of particle~1 through the centre, with the opposite momentum. $a\approx3.6645$
is the equilibrium position.}
\label{tab:cases}
\begin{tabular}{@{}lllll@{}}
\toprule
Case & $\q_1(0)$ & $\p_1(0)$ & $E$ & Motion\\
\midrule
1 & $(a,a)+5\cdot10^{-4}(1,1)$ & $(10^{-3},0)$ & $1.9877$ & small regular oscillation\\
2 & $(a,a)+0.25\,(1,1)$ & $(0,0.05)$ & $2.0073$ & regular\\
3 & $(a,a)+5\cdot10^{-4}(1,1)$ & $(0.1,-0.1)$ & $2.0077$ & close to a periodic loop\\
4 & $(3.8352,\,5)$ & $(0,20)$ & $402.08$ & periodic, vertical bouncing\\
5 & $(3,\,5)$ & $(0,20)$ & $402.00$ & regular, around Case 4\\
\bottomrule
\end{tabular}
\end{table}

\begin{figure}[H]
  \centering
  \includegraphics[width=\textwidth]{fig_cases.pdf}
  \caption{Top: trajectories of the five cases (particle~1 blue, particle~2
  orange). Bottom: 200 returns to a Poincar\'e section, measured from the
  starting point (red star). The section is taken through the starting point,
  across the largest initial velocity component of particle~1: $x_1=x_1(0)$,
  $p_{x_1}>0$ for Cases~1 and~3, so the axes are $(\Delta y_1,\Delta p_{y_1})$;
  $y_1=y_1(0)$, $p_{y_1}>0$ for Cases~2, 4 and~5, so the axes are
  $(\Delta x_1,\Delta p_{x_1})$. Note the scales: in Case~4 all returns lie
  within $10^{-5}$ of the start.}
  \label{fig:cases}
\end{figure}

\subsection{Regular and chaotic motion at increasing energy}\label{sec:sections}
We start from the Case~2 position and give the particles vertical momenta of
increasing size (Figure~\ref{fig:energy-sections}).
\begin{itemize}
  \item At $E\approx2.007$ the section points lie on smooth curves: regular
        motion.
  \item At $E\approx2.028$ and $E\approx2.97$ the points are scattered and the
        trajectory wanders over a large part of the box: chaotic motion.
  \item At $E\approx11$ the motion is regular again. The particles stay in two
        thin vertical strips, close to the stable vertical periodic orbit of
        Section~\ref{sec:vertical}.
\end{itemize}
So for this starting position the motion is chaotic at intermediate energies
and regular again at $E\approx11$, where it stays close to a stable periodic
orbit. Each energy here is a single trajectory, so this does not show that
all motion at $E\approx11$ is regular: Section~\ref{sec:lyapunov} shows that
at high energy chaotic motion is still present, alongside islands of regular
motion.

\begin{figure}[H]
  \centering
  \includegraphics[width=\textwidth]{fig_energy_sections.pdf}
  \caption{Starting from the Case~2 position with vertical momenta $\pm v$ at
  four energies: trajectories (top) and 200 returns to the section
  $y_1=a$, $p_{y_1}>0$ (bottom).}
  \label{fig:energy-sections}
\end{figure}

\subsection{The vertical orbit: shape and stability as the energy changes}\label{sec:vertical}
This orbit is of the type constructed for a rectangular box in
\cite[Sec.~2.6]{turaev2022}: the two particles move fast in the vertical
direction, each on its own line. Because the motion is centre-symmetric,
particle~1 hits the top wall exactly when particle~2 hits the bottom wall.
This is the ``simultaneous impacts'' case, covered by Theorem~5 of
\cite{turaev2022}. That theorem needs the wall potential to behave like
$(\text{distance to the wall})^{-\beta}$ with $\beta>6$, whereas ours has
$\beta=1$, and it holds only at sufficiently high energy. So the theorem
does not guarantee that the orbit is stable in our system; the results below
are a numerical check of whether it is.

Starting from the Case~4 orbit with a smaller momentum ($E\approx11.08$), we
followed the vertical orbit down in energy (Figure~\ref{fig:vertical}).
Newton converged at every energy. The results are:
\begin{itemize}
  \item For $E$ from $11.1$ down to $3.88$, all multipliers have $|\lambda|=1$:
        the orbit is linearly stable.
  \item At $E\approx3.32$, $2.77$ and $2.66$ the orbit is unstable, with largest
        multiplier $|\lambda|\approx2.59$, $6.04$ and $5.21$. So it loses
        stability somewhere between $E\approx3.9$ and $3.3$.
  \item At $E\approx2.55$, Newton no longer found the vertical orbit but a
        different, stable orbit in which the particles move along S-shaped
        paths (Figure~\ref{fig:vertical}(d)).
\end{itemize}
We did not study the transition between $E\approx3.9$ and $3.3$ in detail;
this would need much smaller energy steps.

% Generated by images/make_orbit_figures.py
\begin{figure}[H]
  \centering
  \includegraphics[width=\textwidth]{fig_vertical_family.pdf}
  \caption{One period of the vertical orbit at four energies. (a), (b) stable;
  (c) unstable; (d) at $E\approx2.55$ Newton found a different, S-shaped orbit.
  The black dot is where the orbit crosses the section $y_1=5$.}
  \label{fig:vertical}
\end{figure}

\subsection{The circular orbit: shape and stability as the energy changes}\label{sec:circular}
Newton started directly from Case~3 did not converge, but it did converge from
a nearby symmetric starting point at the same energy, using the diagonal
$x_1=y_1$ as the section. We then followed this orbit upward in energy
(Figure~\ref{fig:circular}). Newton converged at all 21 energies from
$E\approx2.01$ to $4.02$, and every orbit found is linearly stable. As the energy increases, the loop grows towards
the walls. We also computed the orbit at much higher energies, $E\approx52$ and
$E\approx202$: it has become a diamond that touches the middle of each wall,
and it is still linearly stable. The two particles follow the same curve half
a period apart, so this is a choreographic motion in the sense of
\cite{turaev2022}. Diamond-shaped orbits are mentioned there as another family
of periodic motions in a box on which such choreographies could be built, but
they are not analysed \cite[Sec.~6]{turaev2022}; we have not checked how our
orbit is related to that family. Since we only computed the orbit at a finite
set of energies, we cannot rule out short energy intervals in between where it
is unstable.

\begin{figure}[H]
  \centering
  \includegraphics[width=0.8\textwidth]{fig_circular_family.pdf}
  \caption{One period of the circular orbit at several energies. The two
  particles move on the same curve, half a period apart. (a) For
  $2.0\le E\le4.0$ the loop grows towards the walls. (b) At $E\approx52$ and
  $202$ it is a diamond touching all four walls. All orbits shown are linearly
  stable.}
  \label{fig:circular}
\end{figure}

\subsection{Chaotic region and stable islands at \texorpdfstring{$E\approx52$}{E = 52}}\label{sec:lyapunov}
Figure~\ref{fig:lyapunov} shows the Poincar\'e section $y_1=5$, $p_{y_1}>0$
for centre-symmetric motions at $E\approx52.23$, with every point coloured by
the Lyapunov exponent of its trajectory. Most of the section is a chaotic
region with $\lambda_{\max}\approx0.15$--$0.27$. Inside it there are many
islands of regular motion with $\lambda_{\max}\approx0$. The histogram on the
right has two clear peaks, one for regular and one for chaotic trajectories.

At the centre of each island there is a stable periodic orbit.
Figure~\ref{fig:islands} shows ten of them, found with Newton's method. The
simplest one, \#1, is the vertical orbit of Section~\ref{sec:vertical} at this
higher energy: computing the vertical orbit directly at $E\approx52.23$ gives
the same period, $T\approx2.861$, and it is linearly stable.
% Checked 2026-11-08 with continue_to_periodic_orbit: T=2.86096, x1*=3.84204,
% max|lambda|=1.00000001 on the full 6-dimensional section.
The others cross the section two to four times per period and have more
complicated shapes. Because of these islands, a chaotic trajectory cannot
visit the whole section. The theorem of \cite{turaev2022} predicts an island
of this kind only around the vertical orbit, and only under the conditions
discussed in Section~\ref{sec:vertical}; the other islands are not covered by
it.

\begin{figure}[H]
  \centering
  \includegraphics[width=\textwidth]{lyapunov_poincare_section_reduced_report.jpg}
  \caption{Poincar\'e section at $E\approx52.23$, each point coloured by the
  Lyapunov exponent $\lambda_{\max}$ of its trajectory (13\,920 trajectories,
  $T=500$). Dark points ($\lambda_{\max}\approx0$) form islands of regular
  motion in the chaotic region. Right: histogram of $\lambda_{\max}$.}
  \label{fig:lyapunov}
\end{figure}

\begin{figure}[H]
  \centering
  \includegraphics[width=\textwidth]{lyapunov_section_periodic_orbits_report.jpg}
  \caption{Stable periodic orbits at the centres of islands, $E\approx52.23$.
  Top: their points on the section. Below: the orbits in the box, with the
  number $k$ of crossings per period and the period $T$.}
  \label{fig:islands}
\end{figure}

% ---------------------------------------------------------------------------
\section{Discussion}\label{sec:discussion}

\paragraph{Summary.}
Near the equilibrium the motion is regular; slightly above it, for
$E\approx2.03$--$3$, the trajectories we followed are chaotic; at higher
energy chaotic motion is still present, but stable periodic orbits are
surrounded by islands of regular motion. The vertical orbit is stable
at high energy but becomes unstable below $E\approx3.9$. The circular orbit is
stable at every energy we computed, from $E\approx2$ to $E\approx202$.

\paragraph{A difficulty: following the circular orbit downwards.}
Following the circular orbit to energies \emph{below} the Case~3 energy $E_0$
did not work well. With steps of $0.1\%$ in energy, Newton still converged,
but the period of the orbit it found jumped from $101.8$ to $109.2$, $135.4$
and $140.6$, and at $0.998E_0$ it found an unstable, non-symmetric orbit. So
Newton was jumping between different periodic orbits that lie close together,
instead of following one orbit. At $0.996E_0$ Newton did not converge at all;
with half the energy step it converged again. A continuation method that is
designed to stay on one branch, such as pseudo-arclength continuation
\cite{kuznetsov}, would be the natural next step.

\paragraph{Limitations.}
The Lyapunov exponents are computed over a finite time, so very weak chaos
may be missed. The stability results hold at the energies we computed; they
say nothing about short intervals in between. Linear stability also does not
prove that the orbit is stable for larger perturbations.

\paragraph{Open questions.}
\begin{itemize}
  \item What exactly happens to the vertical orbit between $E\approx3.9$ and
        $3.3$, and where does the S-shaped orbit come from?
  \item How does the size of the chaotic region change as the energy grows?
  \item How do the results depend on the exponents $\alpha$, $\beta$ and on the
        shape of the box?
\end{itemize}

% ---------------------------------------------------------------------------
\section*{Code and reproducibility}
All code is available at
\url{https://github.com/jiajunwu1729/2particle-numerics}. Table~\ref{tab:repro} lists the
script that produces each figure.

\begin{table}[H]
\centering
\caption{Figures and the code that produces them.}
\label{tab:repro}
\begin{tabular}{@{}ll@{}}
\toprule
Figure & Source\\
\midrule
\ref{fig:cases}, \ref{fig:energy-sections} & \texttt{images/make\_section\_figures.py} (same computation as \texttt{report03.ipynb})\\
\ref{fig:vertical}, \ref{fig:circular} & \texttt{images/make\_orbit\_figures.py} (periodic points from \texttt{report04.ipynb})\\
\ref{fig:lyapunov} & \texttt{lyapunov\_poincare\_section.py}, then \texttt{images/compress\_figures.py}\\
\ref{fig:islands} & \texttt{lyapunov\_section\_periodic\_orbits.py}, then \texttt{images/compress\_figures.py}\\
\bottomrule
\end{tabular}
\end{table}

\section*{Acknowledgements}
I thank Prof.\ Dmitry Turaev for suggesting this project and for his guidance.

% ---------------------------------------------------------------------------
\begin{thebibliography}{9}
\bibitem{turaev2022}
V.~Rom-Kedar and D.~Turaev,
Stable motions of high energy particles interacting via a repelling potential,
\emph{Comm.\ Math.\ Phys.}\ \textbf{405} (2024), no.~6, 150;
arXiv:2208.14993.
% Verified via Crossref, DOI 10.1007/s00220-024-04946-3 (2026-11-08).

\bibitem{hairer2006}
E.~Hairer, C.~Lubich and G.~Wanner,
\emph{Geometric Numerical Integration}, 2nd ed., Springer, 2006.

\bibitem{lichtenberg1992}
A.~J.~Lichtenberg and M.~A.~Lieberman,
\emph{Regular and Chaotic Dynamics}, 2nd ed., Springer, 1992.

\bibitem{kuznetsov}
Yu.~A.~Kuznetsov, \emph{Elements of Applied Bifurcation Theory}, 3rd ed.,
Springer, 2004.
\end{thebibliography}

\end{document}
